% Author: ChatGPT. Mathematical construction before formalization.
% Input for ChatGPT_ABC_Uniformity_2026.tex; no standalone preamble.
% New bibliography keys: MochCanLift, MochEtaleTheta,
% TakeuchiFuchsian, SijslingCanonical.
\section{Constructing the global initial theta data}
\label{sec:initial-theta-data}

The local fields of Section~\ref{sec:frey-realizations} arise from a
specific rational elliptic curve. We first construct its additional
global covering, cusp, and section of places required by the original
definition of initial theta data, then extend the construction to
every member of Theorem~\ref{thm:powerfree-global-family}.
The construction uses
\cite[Definition~3.1(a)--(f)]{IUT1} directly. It does not invoke an
existence theorem outside a finite exceptional set, and it keeps the
chosen prime and the global base field fixed.

For the first construction put
\begin{equation}\label{eq:itd-curve}
 \begin{gathered}
  \ell=43,\qquad p=1289,\qquad A=p(p^{16}+428),\\
  a=A^2,\qquad b=A^2+1,\qquad c=2A^2+1,\\
  D:\ y^2=x(x-a)(x+b).
 \end{gathered}
\end{equation}
The rational change $x=A^2X$, $y=A^3Y$ identifies $D$ with the
Legendre curve of parameter $-1-A^{-2}$. Define
\begin{equation}\label{eq:itd-fields}
 \begin{gathered}
  F=\Q(i,D[30]),\qquad K=F(D[43])=\Q(i,D[1290]),\\
  X_F=D_F\setminus\{O\},\qquad
  C_F=[X_F/\{\pm1\}],\qquad C_K=C_F\times_FK,\\
  S=\{r:r\text{ is an odd prime and }r\mid abc\}.
 \end{gathered}
\end{equation}
Fix an algebraic closure $\overline F$ containing these fields.
Quotients by inversion are stack quotients; in particular, their
finite etale covering maps are not claims about unramified maps of
coarse curves. The set $S$ is finite and contains 1289.

We first make explicit the two conditions that cannot be replaced by
genericity assertions: stable reduction above 2 and the required core.

\begin{lemma}[The place above 2]\label{lem:itd-two}
Let $k$ be a finite extension of $\Q_2$ and let $E/k$ be an elliptic
curve with rational full $3$-torsion and integral $j$-invariant. Then
$E$ has good reduction over $k$. In particular, $D$ has good reduction
at every place of $F$ above 2.
\end{lemma}

\begin{proof}
Use a minimal Weierstrass equation over the complete valuation ring
of $k$ to define $E_0(k)$ and $E_1(k)$. Its residue field is finite
and perfect, as required in Chapter~VII of \cite{Silverman}.
Suppose that reduction is additive. The group $E_1(k)$ has no
nonzero $3$-torsion, by \cite[VII.3.1(a)]{Silverman}; the quotient
$E_0(k)/E_1(k)$ is the additive group of the residue field, by
\cite[VII.2.1 and VII.5.1(c)]{Silverman}, and also has no
$3$-torsion. Hence $E_0(k)[3]=0$. The rational group
$E[3]\cong(\Z/3\Z)^2$ would inject into $E(k)/E_0(k)$.
The latter group has order at most four in the additive case, by
the Kodaira--Neron bound \cite[VII.6.1]{Silverman}. This is a
contradiction. The cited bound applies in residue characteristic two
and over each finite extension $k$; no unramifiedness assumption is
being made. Finally, multiplicative reduction would force
$v_k(j(E))<0$, whereas the $j$-invariant was assumed integral.
Only good reduction remains.

For \eqref{eq:itd-curve}, $A$ is odd, so $A^2\equiv1\pmod8$.
Writing $S_D=3A^4+3A^2+1$, one has
\[
 j(D)=\frac{256S_D^3}{(abc)^2},\qquad
 v_2(b)=1,\qquad v_2(S_D)=0.
\]
Thus $v_2(j(D))=6$. Every completion of $F$ contains $D[3]$, and
the first assertion applies. Notice that the minimal equation over
the completion need not be the rational equation in
\eqref{eq:itd-curve}.
\end{proof}

The references in this argument have a useful precise location in
the archived second edition of \cite{Silverman}: the Chapter~VII
local-field hypotheses are on printed page~185, and VII.2.1,
VII.3.1(a), VII.5.1(c), and VII.6.1 are on printed pages
188, 192, 196, and 200, respectively. In the archived PDF these are
pages 201, 204, 208, 212, and 216, respectively.

\begin{proposition}[The required core]\label{prop:itd-core}
The hyperbolic orbicurve $C_K$ in \eqref{eq:itd-fields} is a
$K$-core in the sense of \cite[Remark~2.1.1]{MochCanLift}.
The same assertion holds after replacing $K$ by any of its
completions.
\end{proposition}

\begin{proof}
Takeuchi's classification gives exactly four geometric arithmetic
once-punctured elliptic curves
\cite[Theorem~4.1(i)]{TakeuchiFuchsian}. Their $j$-invariants, as
computed in \cite[Section~3.1, Table~4]{SijslingCanonical}, are
\begin{equation}\label{eq:itd-four}
 \frac{2^{14}31^3}{5^3},\qquad
 \frac{2^2 73^3}{3^4},\qquad1728,\qquad0.
\end{equation}
All four are integral at 1289. The curve $D$ satisfies
$v_{1289}(j(D))=-4$, by the split multiplicative calculation in
Section~\ref{sec:frey-realizations}. Thus its once-punctured curve
is not among the four arithmetic curves. Arithmeticity is
unchanged under finite etale coverings, so the hemi-elliptic
quotient is also nonarithmetic over an algebraic closure.

By \cite[Proposition~2.7]{MochCanLift}, a nonarithmetic punctured
hemi-elliptic orbicurve over an algebraically closed field of
characteristic zero is its own core.
Proposition~2.3(i)--(ii) of the same paper gives invariance under
algebraically closed base extension and descent from a separable
closure. Apply these results first over $K$ and then over each
completion. They give the asserted cores.
\end{proof}

The source locations for this step are author-PDF pages 9--10 and
14--15 of \cite{MochCanLift}, published page 392 (PDF12) of
\cite{TakeuchiFuchsian}, and PDF11 of the pinned arXiv version of
\cite{SijslingCanonical}. In particular, the argument does not
replace the four exclusions by a non-CM hypothesis.

\begin{lemma}[Decorated quotient transitivity]\label{lem:itd-decorated}
Let $V$ be a two-dimensional vector space over a field, equipped
with a nondegenerate alternating form. The group
$\operatorname{SL}(V)$ acts transitively on pairs $(H,\bar v)$
consisting of a line $H\subset V$ and a nonzero vector
$\bar v\in V/H$.
\end{lemma}

\begin{proof}
Choose a lift $v$ of $\bar v$. There is a unique $h\in H$ with
$\omega(h,v)=1$. The pair $(h,v)$ is a symplectic basis. Given a
second decorated pair, construct $(h',v')$ in the same way. The
linear map $h\mapsto h'$, $v\mapsto v'$ preserves the alternating
form, has determinant one, and carries the first decorated quotient
to the second.
\end{proof}

\begin{theorem}[Explicit initial theta data]\label{thm:itd-data}
For every nonempty subset $S_0\subseteq S$, there exist a
hyperbolic orbicurve $\underline C_K$, a cusp $\epsilon$ of it,
and a subset $\mathbb V\subseteq V(K)$ mapping bijectively to
$V(\Q)$, such that
\begin{equation}\label{eq:itd-tuple}
 (\overline F/F,X_F,43,\underline C_K,
        \mathbb V,S_0,\epsilon)
\end{equation}
satisfies all of \cite[Definition~3.1(a)--(f)]{IUT1}.
More precisely, one may fix any line $H\subset D[43](K)$ and any
nonzero $\bar v\in D[43](K)/H$ in advance, construct a single
global cover and cusp from them, and then choose $\mathbb V$ with
the required local properties. The fields $F$ and $K$ remain
exactly those in \eqref{eq:itd-fields}.
\end{theorem}

\begin{proof}
We follow the letters of the source definition, whose full statement
is on author-PDF pages 61--63 of \cite{IUT1}.

\smallskip\noindent
\emph{The arithmetic requirements (a)--(c).}
The field $F$ contains $i$ and all $6$-torsion, and is Galois over
$\Q$. Section~\ref{sec:frey-realizations} gives
\begin{equation}\label{eq:itd-degree}
 [F:\Q]\mid
 2|\GL_2(\Z/30\Z)|=276480,
 \qquad 43\nmid[F:\Q],
\end{equation}
and proves that the image of $G_F$ on $D[43]$ contains
$\operatorname{SL}_2(\mathbb F_{43})$. Its kernel cuts out $K$.
Since the pointed curve has a model over $\Q$, every automorphism
of $\overline\Q/\Q$ fixes its geometric isomorphism class.
Consequently its field of moduli is exactly $F_{\rm mod}=\Q$.

Every odd bad prime of $D$ belongs to $S$ and is already split
multiplicative over $\Q$. All other finite primes except possibly
2 are good. Lemma~\ref{lem:itd-two} treats 2. These properties
persist after the extension to $F$. The zero section makes the
smooth genus-one model stable as a one-pointed curve. At a
multiplicative place it gives a stable nodal model, either directly
from the minimal cubic or by contracting the unmarked two-valent
components of the semistable polygon. Thus $X_F$ has stable
reduction at every finite place, in the pointed-curve sense needed
in the definition.

For each $r\in S$, Section~\ref{sec:frey-realizations} proves
\[
 1\le v_r(abc)<43,
 \qquad \operatorname{ord}_w(q_w)=2e(w/r)v_r(abc)
       \quad(w\mid r\text{ in }F).
\]
The ramification index divides \eqref{eq:itd-degree}, so 43 is
prime to all of these positive Tate orders. It is also different
from every $r\in S$, since $(a,b,c)\equiv(1,2,3)\pmod{43}$.
This verifies all the required order and residue-characteristic
conditions for $S_0$, including the choice $S_0=S$. There is no
need to know the complete factorization of the large endpoints:
the finite divisibility certificate and size bound in
Section~\ref{sec:frey-realizations} prove the displayed inequalities.

\smallskip\noindent
\emph{The global covering in (d).}
Put $V=D[43](K)$ and fix $(H,\bar v)$ as in the statement. Take
the quotient isogeny and its dual
\[
 \phi:D_K\longrightarrow D_K/H,
 \qquad \psi:D_K/H\longrightarrow D_K,
 \qquad \psi\phi=[43].
\]
Both are defined over $K$, since $V$ is rational over $K$. Their
kernels give the exact sequence of constant finite groups
\begin{equation}\label{eq:itd-quotient}
 0\longrightarrow H\longrightarrow V
   \xrightarrow{\phi} Q:=\ker\psi\longrightarrow0.
\end{equation}
In particular, $Q\cong V/H$ is cyclic of order 43. Set
\begin{equation}\label{eq:itd-cover}
 \underline X_K=(D_K/H)\setminus Q,
 \qquad \underline C_K=[\underline X_K/\{\pm1\}].
\end{equation}
The restriction of $\psi$ is a degree-43 finite etale cyclic
cover $\underline X_K\to X_K$ with group $Q$. Its cusps are all
$K$-rational. Since the cover extends etale across the origin of
the completed elliptic curve, inertia at that cusp acts trivially
on its fiber. The residue-field part of the cusp decomposition
group also acts trivially, because every point of $Q$ is rational.
The quotient of the fundamental group is therefore trivial on
the full cusp decomposition group and surjective on geometric
monodromy. As an isogeny cover, it factors through the elliptic
abelian quotient. This is precisely the cyclic quotient used in
\cite[Definition~2.1 and its preceding construction]{MochEtaleTheta}.

Inversion yields the cartesian diagram
\begin{equation}\label{eq:itd-square}
 \begin{matrix}
  \underline X_K&\longrightarrow&X_K\\
  \downarrow&&\downarrow\\
  \underline C_K&\longrightarrow&C_K.
 \end{matrix}
\end{equation}
The lower arrow is finite etale by descent along $X_K\to C_K$.
Thus $\underline C_K$ has type $(1,43\text{-tors})^\pm$ in the
source terminology. It is not necessary, and generally not true,
that this lower cover is Galois. By
Proposition~\ref{prop:itd-core}, $C_K$ is a $K$-core. Finite
etale commensurable orbicurves have the same category of finite
etale correspondences, so this is also the $K$-core of
$\underline C_K$.

The nonzero element $\bar v\in V/H$ gives a nonzero cusp in
$Q$ through \eqref{eq:itd-quotient}; its inversion orbit gives a
cusp $\epsilon$ of $\underline C_K$. Both the cover and
$\epsilon$ are now fixed globally. They will not be changed in
the construction at the different primes.

\smallskip\noindent
\emph{The section and the local cusp conditions in (e)--(f).}
For each $r\in S_0$, initially choose any place $w$ of $K$ above
$r$. Split Tate uniformization over $F_{w|F}$ and then $K_w$
gives
\begin{equation}\label{eq:itd-tate-sequence}
 0\longrightarrow\mu_{43}
 \longrightarrow D[43](K_w)
 \longrightarrow\Z/43\Z\longrightarrow0.
\end{equation}
All these torsion points come from $V$. The kernel defines a
line $H_w\subset V$, and the class of $q_w^{1/43}$ determines
a nonzero vector $\bar v_w\in V/H_w$. Changing the root only
multiplies by $\mu_{43}$. Reversing the Tate orientation changes
the vector's sign, which is permitted for the cusp of the
inversion quotient. For the following argument choose one of the
two orientations.

The decorated pairs are equivariant under
$\operatorname{Gal}(K/F)$. With the convention
$(gw)(x)=w(g^{-1}x)$, the map of completions
$g:K_w\to K_{gw}$ transports Tate uniformization and its
multiplicative subgroup. It also transports the graph generator.
Thus
\begin{equation}\label{eq:itd-place-action}
 (H_{gw},\bar v_{gw})=g(H_w,\bar v_w)
\end{equation}
for compatible orientation choices; the equality modulo sign does
not require those choices.
Lemma~\ref{lem:itd-decorated}, together with the actual
$\operatorname{SL}_2(\mathbb F_{43})$ in the Galois image,
therefore gives $g_r\in\operatorname{Gal}(K/F)$ such that
$g_r(H_w,\bar v_w)=(H,\bar v)$. Choose $v_r=g_rw$ as the
selected place above $r$.

The isogeny direction is significant. At this place, write the
curve as $E(q)$ and the matched line as $H=\mu_{43}$. Then
\begin{equation}\label{eq:itd-tate-isogenies}
 \begin{aligned}
  \phi:E(q)&\longrightarrow E(q^{43}),&
           [z]&\longmapsto[z^{43}],\\
  \psi:E(q^{43})&\longrightarrow E(q),&
           [z]&\longmapsto[z].
 \end{aligned}
\end{equation}
The kernel of $\psi$ is generated by $[q]$, and
$\phi([q^{1/43}])=[q]$. Accordingly $\psi$, which defines our
cover, is the graph cover obtained by reducing the natural
$\Pi_X^{\rm tp}\twoheadrightarrow\Z$ quotient modulo 43.
Its distinguished cusp is the image of the generator, up to sign.
This is the graph quotient constructed at the beginning of
\cite[Section~1]{MochEtaleTheta}; its profinite completion is the
$\widehat\Z$ quotient in Definition~2.5(i) of that paper.
It follows that $\underline C_{v_r}$ has type
$(1,\Z/43\Z)^\pm$, and that $\epsilon_{v_r}$ is its
specified $\pm1$ cusp.

At all remaining rational places, including the archimedean place,
choose any one place of $K$ above it. Combining these choices with
the $v_r$ gives $\mathbb V$. The source requires a section of
$V(K)\to V(F_{\rm mod})$, with no equivariance, continuity, or
previously fixed section condition. Thus the $g_r$ can be chosen
independently at the finitely many bad primes. There is still one
global $H$, one global cover, and one global $\epsilon$; the
argument has not substituted different global covers at different
places.

\smallskip\noindent
\emph{The auxiliary theta covers.}
We finish by checking the auxiliary construction recalled in
parts (d)--(f), without extending $F$. For odd $\ell$, inversion
determines the canonical geometric subgroups of the theta cover,
by \cite[Proposition~2.2]{MochEtaleTheta}. An arithmetic splitting
of the indicated cusp extension gives a model over the base
field; the choices form a torsor under
$H^1(G_k,\mu_{43})=k^\times/(k^\times)^{43}$. Part (d) of the
initial-data definition only invokes the geometric subgroups.
Part (e) then invokes their natural local models at each selected
bad place; it does not prescribe a single arithmetic splitting
over $K$ compatible with all local splittings.

Let $k=K_v$ at a selected bad place. Its residue characteristic
is odd and prime to 43, the curve is split multiplicative, and
$i\in k$. The nonarithmeticity needed at the beginning of
\cite[Section~2]{MochEtaleTheta} follows from
Proposition~\ref{prop:itd-core} and its proof. Moreover rational
full 2-torsion implies $q^{1/2}\in k$. Indeed, a square root
$b_0$ represents a rational Tate 2-torsion class. For every
$\sigma\in G_k$, the ratio $\sigma(b_0)/b_0$ belongs to
$q^\Z$ and has valuation zero, so it equals one. Thus
$k=k(\mu_2,q^{1/2})=\ddot k$, the field condition of
\cite[Definition~2.5]{MochEtaleTheta}.
The root-of-unity condition recalled in Remark~1.10.1(ii) of
that paper also holds: rational $D[3]$ and the Weil pairing give
$\mu_3\subset F$, while $i\in F$ gives $\mu_4\subset F$.
Hence $\mu_{12}\subset F\subset k$, without any further field
extension.

The classical theta series in
\cite[Proposition~1.4]{MochEtaleTheta} has value $2i$ plus terms
of positive valuation at the point represented by $i$. This
value is a unit, so a $k$-unit multiple normalizes it to one and
gives a class of standard type in Definition~1.9 of that paper.
Theorem~1.10(iii) there supplies the compatible
$\{\pm1\}$-structure on the cusp torsor. Passing to the
order-43 quotient gives the required splitting class: the two
representatives differing by $-1$ agree in
$k^\times/(k^\times)^{43}$ because $(-1)^{43}=-1$.
Proposition~2.2 with that splitting gives exactly the local
theta cover in Definition~2.5(i). All cusps used here are
$k$-rational; also $\mu_{43}\subset k$ by the Weil pairing on
the rational $43$-torsion.

These are the local models invoked in
\cite[Definition~3.1(e)]{IUT1}. They use no global extension
obtained by adjoining arbitrary $43$rd roots of constants. The
oriented covers at the end of part (f) are then determined by
the fixed triple $(X_K,\underline C_K,\epsilon)$, as in
Definition~1.1 and Remark~1.1.2 of \cite{IUT1}. This verifies
the final auxiliary construction and completes all parts of the
initial-data definition.
\end{proof}

For checking the last part directly, the relevant author-PDF pages
of \cite{MochEtaleTheta} are: the graph quotient on page~11;
$\ddot k=k(\mu_2,q^{1/2})$ on page~16; Proposition~1.4 on
pages~20--21; Definition~1.9 and Theorem~1.10 on pages~27--28;
and Definitions~2.1, 2.5 with Proposition~2.2 on pages~32--36.
The original PDFs and their hashes are recorded separately, so
these locations refer to fixed source objects.

\subsection{A uniform construction along the power-free family}

The covering construction uses no property specific to the number
43 beyond the arithmetic hypotheses already verified. The following
criterion states precisely what can vary.

\begin{proposition}[A rational initial-data criterion]
\label{prop:itd-criterion}
Let $D/\Q$ be an elliptic curve, let $\ell\ge7$ be prime, and set
\[
 F_D=\Q(i,D[30]),\qquad
 K_D=F_D(D[\ell])=\Q(i,D[30\ell]).
\]
Suppose that $D_{F_D}$ has good or multiplicative reduction at every
finite place, that the image of $G_{F_D}$ on $D[\ell]$ contains
$\operatorname{SL}_2(\mathbb F_\ell)$, and that
$v_{p_0}(j(D))<0$ for some rational prime $p_0>5$.
Let $S_0$ be a nonempty finite set of odd rational primes, different
from $\ell$, where $D/\Q$ is split multiplicative. Suppose that
every integral Tate order at a place of $F_D$ above $S_0$ is prime
to $\ell$.

For any line $H\subset D[\ell](K_D)$ and any nonzero
$\bar v\in D[\ell](K_D)/H$, the quotient isogeny and its dual
construct one global cover $\underline C_{K_D}$ and cusp
$\epsilon$. A section $\mathbb V\subset V(K_D)$ of the rational
place map can be chosen such that
\[
 (\overline F_D/F_D,D_{F_D}\setminus\{O\},\ell,
       \underline C_{K_D},\mathbb V,S_0,\epsilon)
\]
satisfies \cite[Definition~3.1(a)--(f)]{IUT1}.
No extension of $F_D$ or $K_D$ is required.
\end{proposition}

\begin{proof}
The universal degree bound
$[F_D:\Q]\mid276480=2^{11}3^3 5$ is prime to $\ell$.
The field is Galois, contains $i$ and rational full 6-torsion,
and the original pointed curve has field of moduli $\Q$.
The mod-$\ell$ kernel cuts out exactly $K_D$.
The reduction hypothesis gives a stable one-pointed model at
every finite place, by the same smooth or nodal construction
used above. The image and Tate-order hypotheses then give all
the other conditions in parts (a)--(c). The orders are imposed
over $F_D$, as in the source, not over $K_D$.

All four invariants in \eqref{eq:itd-four} are integral at every
prime greater than 5. The condition at $p_0$ therefore proves
nonarithmeticity by the same classification and excludes no
additional curves. The proof of Proposition~\ref{prop:itd-core}
gives a core over $K_D$ and over each completion.

Take the quotient $\phi:D\to D/H$ and its dual
$\psi:D/H\to D$ over $K_D$. Their constant kernel sequence is
$0\to H\to D[\ell]\to\ker\psi\to0$. Removing $\ker\psi$
and quotienting by inversion constructs the cover exactly as in
\eqref{eq:itd-cover}. Its cusps are rational, its geometric
monodromy is cyclic of order $\ell$, and its full cusp
decomposition group acts trivially on the fiber at the origin.
It therefore has the type required in part (d), with the same core.
The class $\bar v$ fixes one nonzero global cusp.

Lemma~\ref{lem:itd-decorated} applies over $\mathbb F_\ell$.
For each $r\in S_0$, choose any place $w$ above it and then an
actual $g_r\in\operatorname{Gal}(K_D/F_D)$ carrying its Tate
decorated quotient to the fixed $(H,\bar v)$. Formula
\eqref{eq:itd-place-action} is independent of 43, so selecting
$g_rw$ gives the required local interpretation of that one
global cover and cusp. The dual isogeny is
$E(q^\ell)\to E(q)$, $[z]\mapsto[z]$; its graph generator is
$[q]$, the image of $[q^{1/\ell}]$. All remaining places may
be selected arbitrarily. This verifies the section and generator
requirements in parts (e)--(f).

Finally, the auxiliary-cover argument above only uses oddness of
$\ell$, the original full mod-$\ell$ torsion, and the explicit
local conditions. At a selected place $k$, rational 2-torsion
gives $k=\ddot k$; rational 3-torsion and $i$ give
$\mu_{12}\subset k$; and the theta value at $i$ is a unit
in odd residue characteristic. EtTh's canonical geometric
subgroups and natural local theta splittings therefore apply.
The sign ambiguity disappears modulo $(k^\times)^\ell$ because
$(-1)^\ell=-1$. For the oriented covers, $\ell\ge7$ is
prime to 6, and the original curve's mod-$\ell$ abelian action
is trivial over $K_D$. There is no extra global arithmetic
splitting condition, and no need to adjoin the full
$\ell$-torsion of $D/H$. These are exactly the remaining
assumptions checked in the preceding proof.
\end{proof}

\begin{corollary}[Initial data throughout the unbounded family]
\label{cor:itd-powerfree-family}
For every member $(\ell,p,A)$ of
Theorem~\ref{thm:powerfree-global-family}, put
\[
 \begin{gathered}
  F_A=\Q(i,D_A[30]),\qquad K_A=\Q(i,D_A[30\ell]),\\
  S_A=\{r:r\text{ is an odd prime},\ r\mid A^2(A^2+1)(2A^2+1)\}.
 \end{gathered}
\]
For every nonempty $S_0\subseteq S_A$, every line
$H\subset D_A[\ell](K_A)$, and every nonzero quotient vector
$\bar v$, Proposition~\ref{prop:itd-criterion} supplies initial
theta data with precisely these fields and the chosen prime
$\ell$. In particular $S_0=S_A$ is allowed.
The construction introduces no further largeness threshold on
$\ell$ beyond that in the family existence theorem.
\end{corollary}

\begin{proof}
Every odd bad prime is split multiplicative by the family
calculation, and all other odd primes are good. Odd $A$ gives
$v_2(j(D_A))=6$, so Lemma~\ref{lem:itd-two} gives good reduction
over each completion of $F_A$ at 2. Thus the first hypothesis
of Proposition~\ref{prop:itd-criterion} holds. The family
theorem gives the required $\operatorname{SL}_2$ image over
$F_A$. At its prescribed prime $p\ge30\ell-1>5$,
$v_p(A)=1$ and $v_p(j(D_A))=-4$, proving the core criterion
uniformly. Also $p\in S_A$, so this set is nonempty, while
$(a,b,c)\equiv(1,2,3)\pmod\ell$ gives $\ell\notin S_A$.

For the order condition one must use the precise power-free
statement; it is not asserted that $A^2$ is $\ell$-power-free.
For $w\mid r$ in $F_A$, the three mutually exclusive cases give
\begin{equation}\label{eq:itd-family-orders}
 \operatorname{ord}_w(q_w)=
 \begin{cases}
  4e(w/r)v_r(A),&r\mid A,\\
  2e(w/r)v_r(A^2+1),&r\mid A^2+1,\\
  2e(w/r)v_r(2A^2+1),&r\mid2A^2+1.
 \end{cases}
\end{equation}
Each valuation on the right lies between 1 and $\ell-1$,
and $e(w/r)$ divides the prime-to-$\ell$ degree $[F_A:\Q]$.
Since $\ell$ is odd, all three orders are prime to $\ell$.
This verifies the last hypothesis of the proposition without
the unnecessary stronger inequality $v_r(abc)<\ell$.
The proposition now applies to each chosen $S_0$ and fixed
decorated quotient.
\end{proof}

When $p\in S_0$, changing the selected place by $g_p$ leaves
the exact local field in the family theorem unchanged up to
isomorphism. It still has ramification index $15\ell$, residue
degree two, and total degree $30\ell$. In the notation of
\eqref{eq:pf-uniformizer}, $\beta$ remains a uniformizer and
\[
 v_p(q^{1/(2\ell)})=2/\ell.
\]
Thus the compatible section does not remove the actual Tate unit
or change the valuation to $1/\ell$. The cover and cusp are
single fixed global objects for each family member; there is
no assertion that distinct curves share one covering or one field.

\begin{remark}[Scope of the construction]\label{rem:itd-boundary}
Theorem~\ref{thm:itd-data} and Corollary~\ref{cor:itd-powerfree-family}
verify the original global initial data, including one covering
and cusp with simultaneous local interpretations for each curve.
Neither identifies our prescribed $\ell$ with an existentially
selected prime in \cite[Theorem~5.7.1]{JoshiIV}.
For the isolated level-43 example, no exceptional-set avoidance
is claimed. For the unbounded family, avoidance of a finite set
fixed before the family varies is the separate statement of
Corollary~\ref{cor:powerfree-finite-exception}; it uses the
original compactly bounded domain specified there.

At the selected place above 1289, the completion is still, up
to isomorphism, the exact field of
Section~\ref{sec:frey-realizations}, with ramification index
645, residue degree two, and total degree 1290. Selecting a
Galois-conjugate place changes none of these quantities. The
native root $q^{1/86}$ retains valuation $2/43$.

Initial data alone do not identify the published pilot family
with the explicit native point family, its saturated unit family,
or a principal maximal-order ideal formed before transport.
The exact local pre-ideal/point-hull equality proved separately
applies to its specified native source and integral arrows; this
initial-data construction does not identify them with the whole
published pilot. The standard Bloch--Kato versus $p^{-1}\log$
normalization, the $j^2$-dependent absolute-value rescaling, the
specified order of hull operations, and the remaining Ind3 and
global Frobenius comparisons still need their own compatible proofs.
Neither the initial-data theorems in this section nor the external
core and theta-cover theorems invoked here have been formalized
in Lean in this work.
\end{remark}
